\section{Gauss diagram formulae of low orders}
\label{sec:GDF}

The Gauss diagram formulae used in this paper are listed below.  
$v_2$ uses the formula from \cite{PolyakViro94}. 
$v_3$ uses the formula from \cite{Zhang23}. 
$v_{4,1}$ and $v_{4,2}$ listed here are new expressions. 
% The blue dot is the base point, and the circle is read counterclockwise.  Each black arrow is
% unsigned in the sense of \Cref{sec:Gauss_diagrams}, so an occurrence is
% weighted by the product of the corresponding crossing signs.  In the formula
% for \(J\), a red arrow marked \(+\) is required to match a positive crossing
% and contributes weight \(1\).

\newcommand{\gdfterm}[1]{%
  \raisebox{-.42\height}{%
    \includegraphics[height=1.5cm]{figures/GDF/#1.pdf}}}

\begin{equation*}
  v_2=\gdfterm{v2}.
  \label{eq:gdf-v2}
\end{equation*}
 
% With the normalization \(v_3(3_1^+)=1\), the formula of \cite[Corollary~4.1]{Zhang23} is
\begin{equation*}
  \begin{aligned}
  v_3={}&
  \gdfterm{v3_01}
  +\gdfterm{v3_02}
  +\gdfterm{v3_03}
  +\gdfterm{v3_04}
  +\gdfterm{v3_05}.
  \end{aligned}
  \label{eq:gdf-v3}
\end{equation*}

% The two primitive invariants of order four are given by
\begin{equation*}
  \begin{aligned}
  v_{4,1}={}&
  \gdfterm{J_01}
  +\gdfterm{J_02}
  +\gdfterm{J_03}
  +\gdfterm{J_04}
  +\gdfterm{J_05}
  +\gdfterm{J_06}
  +\gdfterm{J_07}
  +\gdfterm{J_08}
  \\[.4em] 
  &+\gdfterm{J_09}
  +\gdfterm{J_10}
  -\gdfterm{J_11}
  +\gdfterm{J_12}
  +\gdfterm{J_13}
  +\gdfterm{J_14}
  +\gdfterm{J_15}.
  \end{aligned}
  \label{eq:gdf-J}
\end{equation*}

\begin{equation*}
  \begin{aligned}
  v_{4,2}={}&
  2\,\gdfterm{E_01}
  +2\,\gdfterm{E_02}
  +\gdfterm{E_03}
  -\gdfterm{E_04}
  +2\,\gdfterm{E_05}
  +2\,\gdfterm{E_06}
  -\gdfterm{E_07}
  -\gdfterm{E_08}
  \\[.4em]
  &+3\,\gdfterm{E_09}
  +\gdfterm{E_10}
  -\gdfterm{E_11}
  +\gdfterm{E_12}
  +\gdfterm{E_13}
  +3\,\gdfterm{E_14}
  +\gdfterm{E_15}
  +\gdfterm{E_16}
  \\[.4em]
  &+2\,\gdfterm{E_17}
  -\gdfterm{E_18}
  +\gdfterm{E_19}
  +\gdfterm{E_20}
  +\gdfterm{E_21}
  +3\,\gdfterm{E_22}
  +\gdfterm{E_23}
  +\gdfterm{E_24}
  \\[.4em]
  &+\gdfterm{E_25}
  +\gdfterm{E_26}
  +3\,\gdfterm{E_27}
  +\gdfterm{E_28}
  +3\,\gdfterm{E_29}
  -\gdfterm{E_30}
  \\[.4em]
  &+\gdfterm{E_31}
  +\gdfterm{E_32}
  -\gdfterm{E_33}
  +\gdfterm{E_34}.
  \end{aligned}
  \label{eq:gdf-E}
\end{equation*}
